% Generated draft: requires mathematical and source review.
\documentclass[11pt]{article}
\usepackage[margin=25mm]{geometry}
\usepackage{fontspec}
\setmainfont{Latin Modern Roman}
\usepackage{amsmath,amssymb,mathtools,unicode-math}
\setmathfont{Latin Modern Math}
\usepackage{enumitem,needspace,xurl,hyperref}
\setlength{\parindent}{0pt}
\setlength{\parskip}{5pt}
\setlength{\emergencystretch}{3em}
\newcommand{\N}{\mathbb{N}}
\newcommand{\Z}{\mathbb{Z}}
\newcommand{\Q}{\mathbb{Q}}
\newcommand{\R}{\mathbb{R}}
\newcommand{\C}{\mathbb{C}}
\newcommand{\F}{\mathbb{F}}
\newcommand{\A}{\mathbb{A}}
\newcommand{\PP}{\mathbb{P}}
\newcommand{\E}{\mathbb{E}}
\newcommand{\eps}{\varepsilon}
\newcommand{\dd}{\,\mathrm{d}}
\begin{document}
\section*{UnsolvedMath ID: 700352}
\textbf{Source problem code:} catalogue entry
\subsection{Definitions and mathematical statement}
A group \(G\) is soluble if its derived series, defined by \(G^{(0)}=G\) and \(G^{(n+1)}=[G^{(n)},G^{(n)}]\), satisfies \(G^{(d)}=1\) for some integer \(d\geq 0\). Using the standard integral interpretation, \(G\) has type \((FP)_\infty\) if the trivial \(\Z G\)-module \(\Z\) admits a projective resolution
\[
\cdots\longrightarrow P_2\longrightarrow P_1\longrightarrow P_0\longrightarrow\Z\longrightarrow 0
\]
in which every \(P_n\) is finitely generated over \(\Z G\).

\textbf{Question (R. Bieri, 1978).} Is every soluble group \(G\) of type \((FP)_\infty\) constructible in the sense of G. Baumslag and R. Bieri, \emph{Math. Z.} \textbf{151}, no.~3 (1976), 249--257?

The supplied record describes this constructibility condition as the possibility of building \(G\) from the trivial group by finite extensions and HNN extensions. A finite extension of a group \(H\) means a group \(K\) containing an isomorphic copy of \(H\) with \([K:H]<\infty\). An HNN extension of a group \(H\), with associated subgroups \(A,B\leq H\) and an isomorphism \(\varphi:A\to B\), has the form
\[
\left\langle H,t\ \middle|\ t^{-1}at=\varphi(a)\text{ for every }a\in A\right\rangle.
\]
Thus the construction description asks for a finite sequence starting with the trivial group and ending with a group isomorphic to \(G\), each step being a finite extension or an HNN extension of the preceding group. However, the supplied excerpt does not reproduce the cited definition of constructibility or specify restrictions on the HNN extensions, such as an ascending requirement. This construction description is therefore subject to verification against the cited Baumslag--Bieri definition; it is not asserted here to be a verified equivalent formulation.
\subsection{Short English statement}
R. Bieri asks whether every soluble group satisfying the integral homological finiteness condition \((FP)_\infty\) is constructible in the cited Baumslag--Bieri sense. The supplied description uses finite extensions and HNN extensions starting from the trivial group, but the precise permitted construction steps require source verification.
\subsection{Sources}
\begin{itemize}
\item Local UnsolvedMath registry: permanent ID \texttt{700352}, exact source problem code \texttt{catalogue entry}, title \emph{Kourovka Notebook Problem 6.5}. The supplied registry labels the problem open; this status has not been independently verified.
\item \emph{The Kourovka Notebook (Unsolved Problems in Group Theory)}, 6th issue (1978), Problem 6.5, proposed by R. Bieri. The supplied background identifies the consulted edition as the 21st issue edition (2026), \texttt{arXiv:1401.0300v48}. The collection description attributes its editing to E. I. Khukhro and V. D. Mazurov.
\item G. Baumslag and R. Bieri, \emph{Math. Z.} \textbf{151}, no.~3 (1976), 249--257. This is the supplied reference for the meaning of constructible; no article title or full definition is supplied.
\end{itemize}
\end{document}
